\chapter{MARCO REFERENCIAL}

\section{Marco legal y regulaciones}
\hspace{1.25cm}
La viabilidad legal y operativa de las actividades en minería subterránea en el Perú está regida principalmente por el Reglamento de Seguridad y Salud Ocupacional en Minería, aprobado mediante el Decreto Supremo N° 024-2016-EM y sus modificatorias emitidas por el Ministerio de Energía y Minas (MINEM), bajo la supervisión de SUNAFIL y OSINERGMIN~\cite{MINEM_DS024_2016,OSINERGMIN_BoletinGSM_2020}. Esta normativa establece la obligación innegociable de garantizar condiciones ambientales seguras antes y durante el ingreso de los trabajadores a las labores bajo tierra, exigiendo en el artículo 246 el control y registro de contaminantes como monóxido de carbono (CO), dióxido de carbono ($\text{CO}_2$), dióxido de nitrógeno ($\text{NO}_2$) y oxígeno ($\text{O}_2$) por debajo de los Límites de Exposición Ocupacional del Anexo 15~\cite{MINEM_DS024_2016}. Asimismo, los artículos 30, 250 y 254 imponen la clausura de accesos o la paralización inmediata de labores ante fallas en los sistemas de ventilación o concentraciones peligrosas de gases de escape diésel, fundamentando éticamente la justificación formal del proyecto: implementar un sistema robótico cooperativo capaz de realizar verificaciones atmosféricas y trazabilidad de datos en zonas de alto riesgo bajo protocolos \textit{no-entry}, eliminando la exposición directa de los inspectores humanos~\cite{MINEM_DS024_2016,OSINERGMIN_BoletinGSM_2020}.

\section{Identificación de normas y estándares de ingeniería}
\hspace{1.25cm}
La definición de los requisitos de diseño del sistema robótico cooperativo requiere alinearse con estándares técnicos internacionales que garanticen la confiabilidad metrológica, la integridad del hardware y la trazabilidad del dato en entornos subterráneos. En lo respectivo al sensado de atmósferas peligrosas, los Límites de Exposición Ocupacional (OELs) e Indicativos (IOELVs) definidos por la Directiva Europea 2017/164 y las guías técnicas armonizadas exigen que los módulos de monitoreo operen dentro de rangos estrictos de precisión, repetibilidad y tiempo de respuesta para monóxido de carbono (CO), dióxido de carbono ($\text{CO}_2$) y dióxido de nitrógeno ($\text{NO}_2$)~\cite{EU_ExposureDirective_2020,Deveau2015_OEL}. Para asegurar lecturas auditablemente válidas y libres de derivas por interferencias ambientales, la selección e integración de los arreglos de sensores electroquímicos y NDIR debe cumplir con los protocolos de acondicionamiento de señal, filtrado de materia particulada y calibración en campo estipulados en la normativa técnica aplicable para la verificación de mezclas de gas~\cite{Cepa2023_sensors,EU_ExposureDirective_2020}.

\hspace{1.25cm}
Por otro lado, la construcción física y la arquitectura lógica de las plataformas terrestres se rigen por estándares de ingeniería que aseguran la robustez mecatrónica e interoperabilidad del sistema. Las envolventes del instrumental electrónico y los actuadores articulares se dimensionan bajo el código de protección IEC 60529, adoptando un grado mínimo de estanqueidad IP65 para resistir la penetración de humedad relativa elevada y partículas abrasivas en suspensión~\cite{OSINERGMIN_BoletinGSM_2020}. Finalmente, el desarrollo del software de control, el modelado cinemático articulado y la arquitectura de distribución de datos entre el microcontrolador principal y la unidad receptora se estructuran sobre el marco ROS 2 (Robot Operating System) bajo estándares de comunicación DDS (Data Distribution Service) y formalismo de transformaciones espaciales, garantizando una telemetría trazable y sincronizada temporalmente durante la navegación en zonas sin cobertura GNSS~\cite{siciliano2010robotics,KimChoi_2025_MineMapping}.

\section{Identificación de restricciones generales}
\hspace{1.25cm}
El diseño e implementación de los dos robots hexápodos exploradores están delimitados por fronteras físicas, operativas y económicas innegociables impuestas por las condiciones de las labores subterráneas en el Perú~\cite{MINEM_DS024_2016,OSINERGMIN_BoletinGSM_2020}. Desde la perspectiva topográfica y de infraestructura, la geometría de galerías, rampas y piques presenta secciones transversales reducidas y pendientes de hasta doce por ciento ($12\%$). Esto impone restricciones dimensionales rígidas a las plataformas, cuya huella operativa no debe comprometer el libre tránsito en las vías subterráneas, respetando de forma estricta la distancia mínima de un metro ($1.0\text{ m}$) exigida para los pasos peatonales por el Reglamento de Seguridad y Salud Ocupacional en Minería~\cite{MINEM_DS024_2016}. Asimismo, para la unidad accionada por micromotores DC N20, las extremidades articulares de tres grados de libertad (3 DoF) se someten a una restricción cinemática proporcional donde la coxa representa la unidad base ($1x = 3.0\text{ cm}$), el fémur el doble ($2x = 6.0\text{ cm}$) y la tibia el triple ($3x = 9.0\text{ cm}$), minimizando el tamaño de la coxa al límite estructural requerido por el actuador para aliviar el brazo de palanca horizontal y el torque exigido en la base~\cite{InsectInspired_2020}.

\hspace{1.25cm}
Por otro lado, existen restricciones severas vinculadas a la transmisión de datos y la viabilidad presupuestal del prototipado. La atenuación electromagnética provocada por las paredes de roca exige operar mediante procesamiento local y almacenamiento no volátil para misiones desconectadas. En el plano económico, la fabricación de ambos robots hexápodos (tanto el impulsado por motores N20 como el basado en servomotores MG996R) se adecúa a un presupuesto acotado de prototipado universitario. Esto restringe la manufactura al uso de filamento PLA configurado al ocho por ciento ($8\%$) de densidad con infill de rejilla (\textit{grid}) y tres ($3$) capas perimetrales de contorno, parámetro establecido tras verificar que las iteraciones de dos capas resultaban frágiles ante los esfuerzos mecánicos transmitidos por el eje de los motores~\cite{OSINERGMIN_BoletinGSM_2020}.

\section{Compromiso ético y responsabilidad profesional con la seguridad y el bienestar social}
\hspace{1.25cm}
El desarrollo del sistema robótico cooperativo responde a la responsabilidad profesional de aplicar la ingeniería mecatrónica en salvaguarda de la vida y el bienestar del personal en minería subterránea, alineándose con el mandato preventivo de la Ley N° 29783 y el D.S. N° 024-2016-EM al sustituir la inspección humana directa en zonas ciegas por exploradores autónomos que reducen la exposición a atmósferas nocivas y la siniestralidad laboral (ODS 3: Salud y Bienestar)~\cite{MINEM_DS024_2016}. Asimismo, el compromiso ético se manifiesta en garantizar la transparencia metrológica y la trazabilidad de los datos recolectados, promoviendo la democratización tecnológica mediante manufactura aditiva local y optimización de costos para que la mitigación del riesgo operacional sea accesible a la pequeña minería, contribuyendo a un entorno de trabajo seguro y de impacto social positivo~\cite{OSINERGMIN_BoletinGSM_2020}.